\documentclass[11pt,a4paper]{article}

\usepackage[margin=1in]{geometry}
\usepackage{amsmath,amssymb}
\usepackage{bm}
\usepackage{booktabs}
\usepackage{listings}
\usepackage{xcolor}
\usepackage{hyperref}
\usepackage{enumitem}
\usepackage[expansion=false]{microtype}
\usepackage{titlesec}

\hypersetup{
    colorlinks=true,
    linkcolor=blue!50!black,
    urlcolor=blue!50!black,
    citecolor=blue!50!black,
}

\newcommand{\qv}{\mathbf{q}}
\newcommand{\xv}{\mathbf{x}}
\newcommand{\tv}{\mathbf{t}}
\newcommand{\R}{\mathbb{R}}
\newcommand{\norm}[1]{\lVert #1 \rVert}
\DeclareMathOperator*{\argmax}{arg\,max}

\title{\textbf{Norm-Aware Threshold Bounds for Fagin's Threshold Algorithm\\
on Unit-Norm Embeddings}}
\author{DocUVerse --- \texttt{simdq}/\texttt{fagin} engine}
\date{2026-07-13}

\begin{document}
\maketitle

\begin{abstract}
Fagin's Threshold Algorithm (TA) retrieves the exact top-$K$ inner-product
matches by scanning per-dimension sorted lists and halting when the $K$th best
score reaches a frontier threshold $T=\sum_j q_j t_j$. This bound is derived
under a per-dimension (box) constraint only. When the corpus vectors are
$\ell_2$-normalized --- as ours are, since we train and fine-tune
granite-embedding models --- the box bound is badly loose: it corresponds to a
hypothetical document sitting at the frontier in \emph{every} dimension
simultaneously, whose norm is $\gg 1$ and which therefore cannot exist. We
propose replacing the halting bound with the solution to the norm-constrained
maximization $\max\{\,\qv\cdot\xv : q_j(x_j-t_j)\le0\ \forall j,\
\norm{\xv}_2\le 1\,\}$. The frontier face flips direction for $q_j<0$ (the
double-scan walks the list from the opposite end), so the per-dimension box is
written as the signed constraint $q_j(x_j-t_j)\le0$. Reading
the query weights $\qv$ as the gradient $\nabla(\qv\cdot\xv)$, this amounts to
projecting the gradient onto the unit ball rather than sliding its level set out
to the box corner. We give (i) a one-line Cauchy--Schwarz clamp
$T'=\min(\sum_j q_j t_j,\ \norm{\qv}_2)$, and (ii) a tight closed-form
water-filling bound $T^\star$ that provably dominates both. We further show, via
the envelope theorem, that the exact first-order threshold reduction of
advancing list $j$ is $q_j\mu_j$ (its KKT box multiplier), giving a strictly
better sorted-access schedule than the current $q_j\,\Delta t_j$ indicator. All
changes tighten halting only; exactness at $\epsilon=0$ is preserved.
\end{abstract}

\section{Background: the current bound}

The engine (\texttt{fagin\_ta.h}) maintains, for each active dimension
$j$ (those with $q_j\neq0$), a cursor into a list of the corpus values
$Y_{\cdot,j}$ sorted by contribution. After some rounds of sorted access the
\emph{frontier} value in dimension $j$ is $t_j$: every not-yet-seen document
$\xv$ satisfies $x_j\le t_j$ when $q_j>0$, and $x_j\ge t_j$ when $q_j<0$ (the
double-scan walks the list from the opposite end). Both cases are the single
signed statement $q_j(x_j-t_j)\le0$, equivalently $q_j x_j\le q_j t_j$ --- the
per-dimension contribution is bounded by its frontier value regardless of sign.
The classic TA halting
threshold is
\begin{equation}
  T \;=\; \sum_{j} q_j\, t_j,
  \label{eq:classic}
\end{equation}
and TA stops as soon as the $K$th best confirmed score satisfies
$s_{(K)} \ge T-\epsilon$. Correctness follows because for any unseen $\xv$,
\begin{equation}
  \qv\cdot\xv \;=\; \sum_j q_j x_j \;\le\; \sum_j q_j t_j \;=\; T ,
  \label{eq:boxbound}
\end{equation}
where each term is bounded independently. Equality in \eqref{eq:boxbound}
requires $x_j=t_j$ for all $j$ simultaneously --- i.e.\ the document sits at the
frontier corner $\tv=(t_1,\dots,t_D)$.

\section{The unused constraint}

Every corpus vector is $\ell_2$-normalized, $\norm{\xv}_2=1$. The frontier corner
$\tv$ has
\begin{equation}
  \norm{\tv}_2^2 \;=\; \sum_j t_j^2 ,
\end{equation}
which early in the scan (large $t_j$) is on the order of $D$ --- far outside the
unit sphere. The classic bound \eqref{eq:classic} is thus the score of a point
that no document can attain, and it overestimates $T$ substantially, delaying
halting and forcing extra sorted and random accesses.

The tightest valid bound respects \emph{both} the frontier box and the norm:
\begin{equation}
  T^\star \;=\; \max_{\xv\in\R^D}\; \qv\cdot\xv
  \quad\text{s.t.}\quad q_j(x_j-t_j)\le0\ \forall j,\qquad \norm{\xv}_2\le 1 .
  \label{eq:constrained}
\end{equation}

\paragraph{Gradient reading.} With $f(\xv)=\qv\cdot\xv$ we have
$\nabla f=\qv$, constant. Classic TA slides the level set of $f$ outward until it
touches the box corner $\xv=\tv$, ignoring the ball. Problem
\eqref{eq:constrained} instead slides it until tangent to the unit ball, clipped
by whichever frontier faces $x_j\le t_j$ it crosses first. The maximizer is the
gradient direction projected onto the feasible region.

\section{Two bounds}

\subsection{Cheap: the Cauchy--Schwarz clamp}

Dropping the box constraint entirely,
$\max_{\norm{\xv}_2\le1}\qv\cdot\xv=\norm{\qv}_2$ (Cauchy--Schwarz), attained at
$\xv=\qv/\norm{\qv}_2$. Since \eqref{eq:constrained} adds constraints, its value
cannot exceed either relaxation:
\begin{equation}
  T' \;=\; \min\!\Big(\textstyle\sum_j q_j t_j,\ \norm{\qv}_2\Big)
  \;\ge\; T^\star .
  \label{eq:clamp}
\end{equation}
$\norm{\qv}_2$ is a per-query constant, computed once. For unit-norm queries it
caps the early-scan threshold at $1.0$ instead of $\sum_j q_j t_j\gg1$. The
clamp is always valid and essentially free.

\subsection{Tight: water-filling}

The clamp uses \emph{either} the box or the ball; the tight bound uses both. We
solve \eqref{eq:constrained} exactly with Lagrange multipliers. Throughout,
abbreviate the maximum feasible per-dimension contribution
\begin{equation}
  c_j \;:=\; q_j t_j, \qquad a_j^2 \;:=\; q_j^2 .
\end{equation}
Note $c_j\ge0$ for every active dimension: for $q_j>0$ the frontier is
$t_j = \max_{\text{unseen}} x_j$, which is positive early in the scan; for
$q_j<0$ the double-scan makes $t_j$ the \emph{most negative} attainable value, so
$q_j t_j>0$ again. The contribution bound $q_j x_j\le c_j$ is thus a bound by a
nonnegative quantity in both directions.

\paragraph{Lagrangian.} Attach $\lambda\ge0$ to the norm constraint written as
$\tfrac12(\norm{\xv}_2^2-1)\le0$ and $\mu_j\ge0$ to each signed box face
$q_j(x_j-t_j)\le0$. For the maximization,
\begin{equation}
  \mathcal{L}(\xv,\lambda,\bm{\mu})
  \;=\; \sum_j q_j x_j
  \;-\; \lambda\,\tfrac12\!\Big(\sum_j x_j^2 - 1\Big)
  \;-\; \sum_j \mu_j\, q_j (x_j - t_j).
  \label{eq:lagrangian}
\end{equation}

\paragraph{KKT conditions.} At the optimum $(\xv^\star,\lambda,\bm\mu)$:
\begin{align}
  \text{(stationarity)}\quad
    & \partial_{x_j}\mathcal{L}
      = q_j - \lambda x_j - \mu_j q_j = 0
      \;\Longrightarrow\; \lambda x_j = q_j(1-\mu_j),
      \label{eq:stat}\\[2pt]
  \text{(primal feas.)}\quad
    & q_j(x_j-t_j)\le0,\qquad \norm{\xv}_2^2\le1,\\[2pt]
  \text{(dual feas.)}\quad
    & \lambda\ge0,\qquad \mu_j\ge0,\\[2pt]
  \text{(compl.\ slack.)}\quad
    & \lambda\big(\norm{\xv}_2^2-1\big)=0,\qquad \mu_j\,q_j(x_j-t_j)=0.
      \label{eq:cs}
\end{align}

\paragraph{The norm constraint binds.} If $\lambda=0$, \eqref{eq:stat} forces
$q_j(1-\mu_j)=0$; on any dimension with $c_j>0$ the objective still strictly
increases by moving $x_j$ toward its cap, so the unconstrained direction $\qv$
would escape the ball --- contradiction unless the ball is active. Hence
$\lambda>0$ and, by \eqref{eq:cs}, $\norm{\xv^\star}_2=1$.

\paragraph{Two states per dimension.} Complementary slackness $\mu_j q_j(x_j-t_j)
=0$ splits the active dimensions into a capped set $S$ and its complement:

\emph{Interior} ($j\notin S$, box slack so $\mu_j=0$). Then \eqref{eq:stat}
gives
\begin{equation}
  x_j^\star \;=\; \frac{q_j}{2\lambda},
  \label{eq:interior}
\end{equation}
where the effective multiplier is $2\lambda$ once the $\tfrac12$ in
\eqref{eq:lagrangian} is absorbed. This is feasible for the box iff
$q_j x_j^\star \le c_j$, i.e.
\begin{equation}
  \frac{a_j^2}{2\lambda}\;\le\; c_j .
  \label{eq:intfeas}
\end{equation}

\emph{Capped} ($j\in S$, box active so $x_j^\star=t_j$). Then \eqref{eq:stat}
gives the multiplier
\begin{equation}
  \mu_j \;=\; 1 - \frac{2\lambda t_j}{q_j}
        \;=\; 1 - \frac{2\lambda c_j}{a_j^2},
  \label{eq:mu}
\end{equation}
and $\mu_j\ge0$ requires $a_j^2 \le 2\lambda c_j$ --- exactly the negation of
\eqref{eq:intfeas}. The two states partition cleanly at the boundary
$a_j^2 = 2\lambda c_j$.

\paragraph{Sign-safe switch.} Dividing the capping test $a_j^2>2\lambda c_j$ by
$a_j^2=q_j^2>0$ gives $c_j/a_j^2 < 1/(2\lambda)$, i.e.
\begin{equation}
  \rho_j \;:=\; \frac{c_j}{a_j^2} \;=\; \frac{q_j t_j}{q_j^2} \;=\; \frac{t_j}{q_j}
  \;>\;0,\qquad
  j\in S \iff \rho_j < \frac{1}{2\lambda}.
\end{equation}
$\rho_j$ is positive in both scan directions (both $c_j\ge0$ and $a_j^2>0$), so
the ordering is sign-safe; the naive $q_j/t_j$ would flip for $q_j<0$.

\paragraph{Solving for $\lambda$.} Substitute the two states into the active norm
constraint $\sum_j (x_j^\star)^2 = 1$:
\begin{equation}
  \underbrace{\sum_{j\in S} t_j^2}_{\text{capped}}
  \;+\; \underbrace{\sum_{j\notin S}\frac{q_j^2}{4\lambda^2}}_{\text{interior}}
  \;=\; 1
  \;\Longrightarrow\;
  \frac{1}{4\lambda^2}\sum_{j\notin S} a_j^2 \;=\; 1-\sum_{j\in S} t_j^2,
\end{equation}
hence the interior scale factor is the leftover norm budget spread over the
interior mass:
\begin{equation}
  \frac{1}{2\lambda}
  \;=\; \sqrt{\frac{\,1-\sum_{j\in S} t_j^2\,}{\sum_{j\notin S} a_j^2}} .
  \label{eq:lambda}
\end{equation}

\paragraph{Back-substitution.} The optimal value splits over the two sets:
\begin{align}
  T^\star
   &= \sum_{j\in S} q_j x_j^\star + \sum_{j\notin S} q_j x_j^\star
    = \sum_{j\in S} q_j t_j + \sum_{j\notin S} \frac{q_j^2}{2\lambda}
      \notag\\
   &= \sum_{j\in S} q_j t_j
    + \frac{1}{2\lambda}\sum_{j\notin S} a_j^2
    = \sum_{j\in S} q_j t_j
    + \sqrt{\frac{1-\sum_{j\in S}t_j^2}{\sum_{j\notin S}a_j^2}}\;\sum_{j\notin S} a_j^2,
\end{align}
using \eqref{eq:interior} then \eqref{eq:lambda}. Collapsing the last product
yields the closed form
\begin{equation}
  \boxed{\;
  T^\star \;=\; \sum_{j\in S} q_j t_j
  \;+\; \sqrt{\Big(\sum_{j\notin S} q_j^2\Big)\Big(1-\sum_{j\in S} t_j^2\Big)}
  \;}
  \label{eq:waterfill}
\end{equation}
The first term is the capped dims pinned at their frontier; the second is
Cauchy--Schwarz applied to the interior dims sharing the leftover norm budget
$\sqrt{1-\sum_{j\in S}t_j^2}$ along $\qv$.

\paragraph{Dual form (what the kernel evaluates).} Rather than search for the
prefix $S$ directly, the implementation dualizes only the norm constraint and
keeps the box explicit, giving the one-dimensional convex dual
\begin{equation}
  g(\lambda) \;=\; \lambda
    \;+\; \sum_j \max_{\,q_j x_j \le c_j}\big(q_j x_j - \lambda x_j^2\big),
  \qquad
  T^\star = \min_{\lambda\ge0} g(\lambda).
  \label{eq:dual}
\end{equation}
Each inner maximization is a concave parabola in $x_j$ peaking at
$x_j=q_j/(2\lambda)$ with unconstrained value $a_j^2/(4\lambda)$; if that peak
respects the cap ($a_j^2\le2\lambda c_j$) it is attained, otherwise the maximum
sits at the boundary $x_j=t_j$ with value $c_j-\lambda t_j^2$:
\begin{equation}
  \max_{q_j x_j\le c_j}(q_j x_j-\lambda x_j^2)
   =\begin{cases}
      \dfrac{a_j^2}{4\lambda}, & a_j^2\le 2\lambda c_j\ (\text{interior}),\\[8pt]
      c_j-\lambda t_j^2,       & \text{otherwise}\ (\text{capped}).
    \end{cases}
  \label{eq:dualdim}
\end{equation}
Because $g$ is convex in $\lambda$, a ternary search on $[\,0,\ \norm{\qv}_2\,]$
converges to $T^\star$; the recovered $\lambda$ feeds $\mu_j$ via \eqref{eq:mu}
for the schedule of Section~4. Crucially, weak duality gives $g(\lambda)\ge
T^\star$ for \emph{every} $\lambda\ge0$, so numerical error in the search can only
\emph{loosen} the bound --- never drop it below $T^\star$ and break the halting
guarantee.

\paragraph{Finding $S$.} Sort dimensions by the sign-safe ratio
$\rho_j=t_j/q_j$ \emph{ascending}; the capped set $S$ is the prefix of smallest
$\rho_j$. Sweep the prefix boundary and accept the first split consistent with
$\rho_j<1/(2\lambda)$ for $j\in S$ and $\ge1/(2\lambda)$ otherwise, where
$1/(2\lambda)=\sqrt{(1-\sum_{j\in S}t_j^2)/\sum_{j\notin S}q_j^2}$ is recomputed
along the sweep. Cost $O(D\log D)$, or $O(D)$ with a selection-based search.
Because the max contribution $q_j t_j$ only \emph{decreases} as the scan proceeds
(so $\rho_j=q_j t_j/q_j^2$ decreases, regardless of $\mathrm{sign}(q_j)$), dims
migrate into $S$ monotonically --- $S$ \emph{grows} over the scan --- enabling
incremental updates instead of a per-round re-sort.

\paragraph{Ordering.} \eqref{eq:waterfill} is provably
$\le\min(\sum_j q_j t_j,\ \norm{\qv}_2)$ --- it is never worse than classic
\emph{or} the clamp --- and it interpolates smoothly: early scan (ball binds,
$S=\varnothing$) gives $T^\star\approx\norm{\qv}_2$; late scan
($\norm{\tv}<1$, box binds, $S=$ all) recovers $T^\star=\sum_j q_j t_j$ and hence
the exact classic bound.

\paragraph{Worked example.} $\qv=(0.6,0.8)$ (so $\norm{\qv}_2=1$),
frontier $\tv=(0.5,0.95)$. Classic \eqref{eq:classic}:
$T=0.6(0.5)+0.8(0.95)=1.06$. Clamp \eqref{eq:clamp}: $\min(1.06,1.0)=1.0$.
Water-filling \eqref{eq:waterfill}: dim~1 caps (smaller ratio,
$\rho_1=0.5/0.6=0.833<\rho_2=0.95/0.8=1.19$),
$T^\star=0.6(0.5)+\sqrt{0.8^2\,(1-0.5^2)}=0.3+\sqrt{0.64\cdot0.75}=0.993$.
Tighter than both, still a valid upper bound.

\paragraph{Negative weights and skipped dims.} The signed constraint
$q_j(x_j-t_j)\le0$ makes the derivation sign-agnostic: for $q_j<0$ the
double-scan frontier is $x_j\ge t_j$, but the contribution bound $q_j x_j\le q_j
t_j$, the interior point $x_j=q_j/(2\lambda)$, and the closed form
\eqref{eq:waterfill} (which uses only $q_j^2$ and $q_j t_j$) all carry over
verbatim. Only the ratio must be the sign-safe $\rho_j=t_j/q_j$, never $q_j/t_j$.
Dimensions with $q_j=0$ are excluded from the active set; bounding only
$\norm{\xv_{\text{active}}}_2\le1$ is conservative and stays valid.

\section{Sharper sorted-access schedule}

The \texttt{steepest} schedule currently keys its dimension heap on the marginal
contribution drop $\delta_j=q_j\,\Delta t_j = c_j(d)-c_j(d{+}\text{batch})\ge0$,
where $c_j=q_j t_j$ (Quick-Combine indicator; sign-safe as written). Under the
classic bound $T=\sum_j c_j$, $\partial T/\partial c_j=1$, so the drop in $T$ is
exactly $\delta_j$. Under the tight bound, the envelope theorem instead gives
\begin{equation}
  \frac{\partial T^\star}{\partial c_j} \;=\; \mu_j \;=\;
  \begin{cases}
    1-2\lambda\,t_j/q_j, & j\in S \ (\text{capped}),\\[2pt]
    0, & j\notin S \ (\text{interior}),
  \end{cases}
  \qquad
  \frac{\partial T^\star}{\partial t_j}=q_j\mu_j,
  \label{eq:envelope}
\end{equation}
with $\mu_j\in[0,1]$. \textbf{Advancing an interior (non-binding) list does not
lower the threshold at all} ($\mu_j=0$). The exact per-round reduction in
$T^\star$ is $\mu_j\,\delta_j$, nonzero only on binding dims --- a strictly better
indicator than $\delta_j$ alone, which spends sorted accesses on dims that are
not moving $T^\star$. The fix is one factor: multiply the existing heap key by
$\mu_j$. This is sign-safe (both $\delta_j\ge0$ and $\mu_j\ge0$ regardless of
$\mathrm{sign}(q_j)$) and directs each sorted access to the list that actually
pulls the threshold down fastest.

\section{Implementation (as built)}

The kernel hooks are in \texttt{docuverse/engines/retrieval/simdq/\_native/%
include/fagin\_ta.h}; the new bound is exposed as two additional
\emph{schedules} rather than a flag, so it flows through the existing
\texttt{schedule} selector end to end.

\paragraph{General ball radius.} The corpus need not be exactly unit-norm, and
assuming so is unsafe: on a non-normalized corpus a $\norm{\qv}_2$ cap is invalid
and silently breaks exactness. We therefore constrain $\norm{\xv}_2\le R$, where
$R^2$ is the \emph{maximum squared row norm}, measured once at build time
(\texttt{FaginIndex.max\_norm\_sq}, persisted in \texttt{meta.json}) and passed
to the kernel. The bound \eqref{eq:waterfill}/\eqref{eq:clamp} generalizes by
$1\!\to\!R^2$ under the norm budget and $\norm{\qv}_2\!\to\!R\norm{\qv}_2$ in the
clamp; it is exact for any $R$ at least the true maximum, and tightest ($R=1$)
for a normalized corpus.

\paragraph{Dual solve.} Rather than the prefix water-filling of
Section~4.2, the kernel minimizes the Lagrangian dual
$g(\lambda)=\lambda R^2+\sum_j\max_{q_j x_j\le q_j t_j}(q_j x_j-\lambda x_j^2)$
over $\lambda\ge0$ by ternary search ($g$ is convex in $\lambda$), then clamps by
$\min(\sum_j q_j t_j,\ R\norm{\qv}_2)$. This yields the same optimum
\eqref{eq:waterfill} but is branch-light and, crucially, \emph{robust}: by weak
duality $g(\lambda)\ge T^\star$ for \emph{every} $\lambda\ge0$, so numerical
error in the $\lambda$ search can only loosen the bound, never drop it below
$T^\star$ and violate correctness. The KKT box multiplier for the schedule key
\eqref{eq:envelope} is read off the converged $\lambda$ as
$\mu_j=\max(0,\,1-2\lambda t_j/q_j)$.

\begin{enumerate}[nosep]
  \item \texttt{fagin\_threshold\_norm(c, a2, t2, ndims, sum\_c, qnorm, r2,
        lam)} --- the dual routine returning $T^\star$ and the optimal $\lambda$.
  \item Two new schedule enums,
        \texttt{FAGIN\_SCHEDULE\_LOCKSTEP\_NORM = 2} and
        \texttt{\_STEEPEST\_NORM = 3}, pair the norm-aware halting bound with the
        existing round-robin and steepest access patterns. \texttt{STEEPEST\_NORM}
        additionally scales the heap key $\delta_j$ by $\mu_j$
        \eqref{eq:envelope}.
  \item Python surface: \texttt{FaginIndex.SCHEDULES} gains
        \texttt{lockstep\_norm}/\texttt{steepest\_norm}; the experiment harness
        prints these as \textbf{GTA}/\textbf{GTASD} (vs.\ TA/TASD for the classic
        box bound).
\end{enumerate}

\paragraph{Correctness \& measured effect.} The change only tightens an upper
bound used for halting, so the returned top-$K$ is identical --- verified byte
for byte against the classic schedules and against brute force
(\texttt{tests/test\_fagin\_*.py}, all schedules parametrized). Exactness holds
at recall@10 $=1.0$ across unit-norm, non-normalized ($R^2\approx65$), and
varying-norm ($R^2\approx530$) corpora. On sorted-access counts (exact,
$\epsilon=0$, $N=3000$, $D=32$): GTASD cuts sorted accesses by $\sim51\%$ on a
unit-norm corpus, $\sim24\%$ non-normalized, and $\sim19\%$ varying-norm --- the
gain grows as the corpus concentrates toward the sphere, where the box bound
overestimates most. GTA (lockstep) matches TA at high $D$ (lockstep halts deep
enough that $\norm{\tv}<R$ and the bound collapses onto the box) but saves up to
$\sim39\%$ in the low-$D$, shallow-halt regime.

\end{document}
